\chapter{Sistema de Autenticación y Seguridad}

\section{Autenticación de usuarios}
La seguridad en el acceso al sistema es un pilar fundamental para proteger la información académica y los datos personales de los usuarios y solicitantes. Se ha implementado un esquema de autenticación híbrido que combina el método tradicional con tecnologías modernas de Single Sign-On (SSO).

\begin{itemize}
    \item \textbf{Autenticación Base (Laravel UI):} Se utiliza el andamiaje proporcionado por \texttt{laravel/ui} para gestionar el registro, inicio de sesión y recuperación de contraseñas. Todas las credenciales son almacenadas en la base de datos utilizando algoritmos de hashing seguros (Bcrypt), garantizando que las contraseñas nunca sean legibles en texto plano.
    \item \textbf{Integración con Google (Laravel Socialite):} Para agilizar el acceso de la comunidad universitaria y los visitantes, se integró el paquete \textbf{Laravel Socialite}. Esto permite iniciar sesión utilizando cuentas de Google (particulares o institucionales \texttt{@unsa.edu.ar}). La tabla \texttt{users} vincula de forma segura el identificador único \texttt{google\_id} con el perfil local del usuario.
    \item \textbf{Redirección Condicional por Perfil:} En \texttt{LoginController} y \texttt{HomeController}, se implementa una lógica de enrutamiento que redirige de forma transparente a los usuarios con rol regular (\texttt{USUARIO}) hacia su panel personal de autogestión (\texttt{route('mis-turnos')}), restringiendo el acceso al panel administrativo (\texttt{/home}) exclusivamente a operadores y administradores.
\end{itemize}

\section{Gestión de roles y permisos}
El control de acceso basado en roles (Role-Based Access Control) se ha implementado de manera profunda utilizando el paquete \textbf{Spatie Laravel-Permission}. Este esquema permite una separación estricta de funciones:

\begin{itemize}
    \item \textbf{Super Administrador / Administrador General:} Posee acceso irrestricto a todos los módulos del sistema, incluyendo la parametrización de dependencias, gestión de roles y permisos, auditoría de usuarios y la gestión exclusiva de plantillas de correo electrónico transaccionales (\texttt{PlantillasEmailController}).
    \item \textbf{Administrador de Dependencia:} Tiene facultades para configurar los trámites, mesas habilitadas, interruptores de servicio y el calendario de feriados correspondientes a su ámbito institucional (ej. el Observatorio Alanis).
    \item \textbf{Operador de Mostrador y Atención:} Rol asignado al personal que atiende al público presencial. Sus privilegios abarcan la visualización de turnos de la jornada, cambio de estados (Atendido, Ausente), llamado secuencial de turnos en cola (\texttt{llamar\_siguiente}) y emisión de pases y derivaciones internas.
    \item \textbf{Usuario Regular / Visitante:} Rol predeterminado para el público general, con permisos acotados exclusivamente a su bandeja personal de autogestión (\texttt{mis-turnos}), descarga de comprobantes propios y cancelación anticipada de sus citas.
\end{itemize}

La protección se aplica tanto a nivel de \textbf{rutas} mediante middlewares agrupados (ej. \texttt{middleware(['role:ADMINISTRADOR|OPERADOR'])}) como a nivel de \textbf{vistas} mediante directivas Blade (ej. \texttt{@can('gestionar\_tramites')}), adaptando la interfaz dinámicamente a los privilegios del usuario autenticado. A nivel de datos, la relación \texttt{usuarios\_dependencias} garantiza la segregación departamental de la información.

\section{Protección de rutas y datos}
El sistema incorpora múltiples capas de protección nativas de Laravel para mitigar las vulnerabilidades web más comunes:

\begin{itemize}
    \item \textbf{Prevención de CSRF (Cross-Site Request Forgery):} Todas las peticiones mutables (POST, PUT, DELETE) requieren un token criptográfico generado por la sesión, invalidando cualquier intento de falsificación de peticiones en sitios cruzados.
    \item \textbf{Validación mediante FormRequests:} Se utilizan clases de validación específicas para sanear y estructurar los datos recibidos (validación de DNI argentino, correos electrónicos válidos, cupos y formatos de fecha) antes de ingresar a la lógica de negocio.
    \item \textbf{Protección contra Inyección SQL:} El uso de \textbf{Eloquent ORM} y el constructor de consultas parametrizadas de PDO garantiza el escape automático de variables, neutralizando ataques de inyección SQL.
    \item \textbf{Acceso Público Seguro por Código de Reserva:} En los flujos públicos de carga de integrantes de delegaciones escolares (\texttt{reserva/\{codigo\}/integrantes}) y cancelación directa (\texttt{reserva/\{codigo\}/cancelar}), la seguridad se basa en la no divulgación del código alfanumérico único generado algorítmicamente, permitiendo a los directivos operar sin crear cuentas previas pero impidiendo el acceso a reservas ajenas.
    \item \textbf{Seguridad en el Lado del Cliente:} Los componentes de \textbf{Vue.js 2} se comunican exclusivamente mediante endpoints protegidos que validan la sesión y los permisos en el servidor, evitando la exposición de reglas críticas de negocio en el navegador.
\end{itemize}
